\documentclass[11pt,a4paper]{article}

\usepackage[margin=1in]{geometry}
\usepackage{amsmath,amssymb}
\usepackage{booktabs}
\usepackage{longtable}
\usepackage{tabularx}
\usepackage{array}
\usepackage{enumitem}
\usepackage{hyperref}
\usepackage{microtype}
\usepackage{xcolor}
\usepackage{listings}
\usepackage{fancyhdr}
\usepackage{titlesec}

\hypersetup{
  colorlinks=true,
  linkcolor=blue!50!black,
  citecolor=blue!50!black,
  urlcolor=blue!50!black
}

\newcolumntype{Y}{>{\raggedright\arraybackslash}X}
\newcommand{\layer}[1]{\paragraph{#1.}}

\setlist[itemize]{leftmargin=*}
\setlist[enumerate]{leftmargin=*}

\lstset{
  basicstyle=\ttfamily\small,
  frame=single,
  breaklines=true,
  columns=fullflexible,
  aboveskip=0.8em,
  belowskip=0.8em
}

\title{Large Language Models --- Concepts From Scratch\\
\large A detailed phase-by-phase paper for reading and interview prep}
\author{Narration: Setup $\rightarrow$ Problem $\rightarrow$ Root cause $\rightarrow$ Solution $\rightarrow$ Payoff\\[0.4em]
Style source:
\href{https://www.youtube.com/watch?v=-AB6m0Spo6c}{Vizuara --- LLM Interview Series no.~5 (PagedAttention)}\\[0.4em]
Do not invent benchmarks. Numbers from papers/blogs are labeled as that source's setting.}
\date{}

\begin{document}
\maketitle

\begin{abstract}
This paper walks the LLM stack in order: what an LLM is, mental models, Transformer and attention, pretraining, post-training, RAG versus fine-tuning, then inference systems (KV cache, GQA, FlashAttention, PagedAttention, quantization, speculative decoding).
Every major idea is told as a story---the naive approach, why it breaks, the root cause, then the method---not as a glossary of acronyms.
Claims stay with checkable papers, the public vLLM blog, and the cited Vizuara teaching videos.
Companion lecture notes: \texttt{research/llm/vizuara\_inference\_phase1\_notes.md} and \texttt{research/llm/llm\_end\_to\_end\_survey.md}.
Markdown twin: \texttt{research/llm\_concepts\_from\_scratch.md}.
\end{abstract}

\tableofcontents
\clearpage

\clearpage
\section{Phase 0 --- The map before the details}


Read this once so later sections attach to a spine.


\begin{lstlisting}
World text
  $\rightarrow$ tokenize into subwords
  $\rightarrow$ pretrain a decoder-only Transformer (base model = next-token predictor)
  $\rightarrow$ post-train (SFT $\rightarrow$ preferences / reasoning)
  $\rightarrow$ product choice: RAG and/or fine-tuning
  $\rightarrow$ serve autoregressive decode:
        prefill vs decode
        $\rightarrow$ KV cache (good compute / evil memory)
        $\rightarrow$ shrink KV (MQA / GQA / \ldots{})
        $\rightarrow$ FlashAttention (IO schedule)
        $\rightarrow$ continuous batching + PagedAttention (multi-tenant memory)
        $\rightarrow$ quantization $\rightarrow$ speculative decoding
  $\rightarrow$ evaluate capability + preference + safety + cost
\end{lstlisting}


\begin{table}[ht]
\centering\small
\begin{tabularx}{\textwidth}{@{}YY@{}}
\toprule
\textbf{Symptom} & \textbf{Likely stage} \\
\midrule
Cannot do basics even with good prompts & Pretraining data / capacity \\
Knows facts, ignores instructions & SFT \\
Fluent but rude / unsafe / sycophantic & Preference post-training \\
Chat OK, hard math fails & Reasoning post-training or test-time compute \\
Correct but too slow / expensive & Inference systems / routing \\
Leaderboard up, users unhappy & Evaluation mismatch \\
\bottomrule
\end{tabularx}
\end{table}


\bigskip


\clearpage
\section{Phase 1 --- What an LLM actually is}


\subsection{Step 1.1 --- Stop saying ``understands language''}


\paragraph{Setup.}

People say ChatGPT ``understands'' Spanish travel plans. Under the hood, the deployed object is a parameterized function: given tokens so far, it outputs a probability distribution over the next token in a fixed vocabulary.


\paragraph{Problem.}

Language is open-ended. Hand-written grammar rules and templates do not cover web text, code, dialogue, and instructions at scale.


\paragraph{Root cause.}

We need one training objective that (i) uses abundant unlabeled text and (ii) produces a generative system.


\subsubsection{Solution --- Autoregressive language modeling}

Factorize a sequence with the chain rule:


\[
P(x_1,\ldots,x_T)=\prod_{t=1}^{T} P(x_t \mid x_{<t}).
\]


Train by minimizing next-token \textbf{cross-entropy} on large corpora. At inference, choose the next token by sampling from that distribution (or taking \(\arg\max\)), append it, and repeat.


\paragraph{Payoff.}

One scalable objective can absorb many downstream behaviors as ``continue this prompt usefully.'' This is the standard framing used across GPT-style models.


\textbf{Anchors:} Brown et al., GPT-3 (2020); standard LM factorization. We are not claiming a specific chatbot's private training recipe.


\subsection{Step 1.2 --- Tokenization (why text is not raw characters)}


\paragraph{Setup.}

Neural nets need integer IDs. Text must become a sequence of token IDs.


\paragraph{Problem.}

\begin{itemize}
\item \textbf{Character-level:} sequences become very long $\rightarrow$ attention cost explodes.  
\item \textbf{Word-level:} vocabulary explodes; rare words, typos, and code break.
\end{itemize}


\paragraph{Root cause.}

We need open vocabulary without pathological sequence length.


\paragraph{Solution.}

\textbf{Subword} tokenization. Common in GPT-style models: Byte-Pair Encoding (BPE) / byte-level BPE --- merge frequent symbol pairs into subwords. Other families use WordPiece or Unigram/SentencePiece.


\paragraph{Payoff.}

Any Unicode string can be encoded; frequent words stay short; rare pieces split.


\bigskip


\clearpage
\section{Phase 2 --- Mental models you will reuse later}


These are not serving kernels. They are shared intuitions under loss, architecture choices, and sampling. Learn them before optimizing decode.


\subsection{Step 2.1 --- Entropy (surprise)}


\textbf{Video:} \href{https://www.youtube.com/watch?v=q3agYqgqklU}{Entropy: AI Mental Model \#9}   \textbf{Visual notes:} \href{https://vizuaraai.github.io/great-mental-models-of-ai/lecture-09-entropy.html}{lecture-09-entropy.html}


\subsubsection{How to explain it (story first --- do not open with \(-\sum p\log p\))}


\begin{enumerate}
\item Everyday disorder (messy room, cream in coffee).  
\item Formal meaning: \textbf{count of arrangements} that look the same from outside.  
\item Physics: heat = molecular jiggling $\rightarrow$ more microstates $\rightarrow$ higher entropy; Second Law.  
\item Boltzmann: \(S = k \log W\).  
\item Shannon: information = \textbf{surprise} \(=-\log p\); entropy = \textbf{expected surprise}.  
\item Only then: where LLMs use the same idea.
\end{enumerate}


\textbf{Core reframing:} Entropy is not a vibe of ``messiness.'' Low entropy $\approx$ few arrangements (tidy / peaked / certain). High entropy $\approx$ many arrangements (messy / flat / uncertain).


\subsubsection{Shannon side}


\paragraph{Setup.}

Shannon (1948) needed to measure how much information a message carries.


\paragraph{Problem.}

``The sun rose'' and ``it snowed in the desert'' are not equally informative.


\paragraph{Root cause.}

Common events are unsurprising; rare events are surprising.


\paragraph{Solution.}

\[
\text{surprise}(x) = -\log p(x)
\]

Entropy of a distribution = average surprise. Fair coin $\rightarrow$ high uncertainty. Biased coin $\rightarrow$ lower average surprise.


\subsubsection{Where it shows up in LLM land}


\textbf{(1) Cross-entropy loss = training surprise.}   If the true next token has model probability \(p_{\text{true}}\):

\[
\mathcal{L} = -\log p_{\text{true}}.
\]

Pretraining is largely minimizing this next-token cross-entropy.


\textbf{(2) Perplexity.}   Perplexity $\approx$ effective branching factor of the model's uncertainty (teaching English: perplexity 2 $\approx$ coin-flip unsure; perplexity 50 $\approx$ \textasciitilde{}50 doors). Lower perplexity $\Rightarrow$ more confident next-token predictions \textit{on that data}. It is an intrinsic LM metric, not the same as chat quality.


\textbf{(3) Sampling temperature.}   Softmax with temperature \(T\): low \(T\) $\rightarrow$ peaked (low entropy, deterministic); high \(T\) $\rightarrow$ flatter (high entropy, more random). The knob is called temperature because it spreads probability mass the way physical temperature spreads molecular arrangements.


\textbf{(4) KL divergence = relative entropy.}   Extra surprise from using the wrong distribution. Shows up as the \textbf{KL penalty} in RLHF/PPO that keeps the policy near a reference model.


\textbf{Interview one-liner:} Whenever the problem is uncertainty, surprise, confidence, randomness, or distance between beliefs --- reach for entropy. In LLM land: CE trains, perplexity scores uncertainty, temperature spends entropy at decode, KL limits drift in RLHF.


\subsection{Step 2.2 --- Expressivity (lift when stuck)}


\textbf{Video:} \href{https://www.youtube.com/watch?v=YE0udJCgDqc}{Expressivity: AI Mental Model \#6}


\paragraph{Story.}

If patterns cannot be separated in the current representation (classic XOR picture), \textbf{lift} to a richer / higher-dimensional space where they become separable (kernel trick, wider FFN, multi-head perspectives).


\paragraph{Warning.}

More expressivity can also fit noise $\rightarrow$ overfitting. In Transformers, the \textbf{FFN} is the per-token ``wide lift''; attention is the inter-token mixer. Do not confuse those jobs.


\subsection{Step 2.3 --- Compression (narrow waist)}


\textbf{Video:} \href{https://www.youtube.com/watch?v=hGS6RbAYLl0}{Compression: AI Mental Model \#5}


\paragraph{Story.}

When something is too big (memory, \(\Delta W\), pixels), find a \textbf{narrow latent}, do the work there, project back --- because the giant object was never as high-dimensional as it looked.


\subsubsection{Same gesture in different places}

\begin{table}[ht]
\centering\small
\begin{tabularx}{\textwidth}{@{}YYY@{}}
\toprule
\textbf{Place} & \textbf{Too big} & \textbf{Waist} \\
\midrule
Autoencoder & Input dim & Latent code \\
LoRA & Full \(\Delta W\) & \(A\,(D\times r)\), \(B\,(r\times D)\), \(r\ll D\) \\
GQA / MLA-style & Full MHA KV & Shared / latent KV \\
Embedding stores & Float embeddings & Quantized / compressed vectors \\
\bottomrule
\end{tabularx}
\end{table}


\textbf{Do not confuse:} GQA changes attention KV layout at architecture/serve time; LoRA compresses fine-tuning updates.


\subsection{Step 2.4 --- Diffusion (contrast with autoregressive LLMs)}


\textbf{Video:} \href{https://www.youtube.com/watch?v=bTEfdB2D1Ek}{Diffusion: AI Mental Model \#7}


\subsubsection{Story (not DDPM jargon first)}

\begin{enumerate}
\item Creating a masterpiece in one shot is hard.  
\item Trick: \textbf{learn to destroy gently}, then run destruction backward.  
\item Forward: add a little noise repeatedly until nearly pure noise (and know what was added).  
\item Train a network to undo \textbf{one tiny step}.  
\item Generate: start from pure noise; repeatedly remove predicted noise.
\end{enumerate}


\subsubsection{Contrast with ChatGPT-class inference}


\begin{table}[ht]
\centering\small
\begin{tabularx}{\textwidth}{@{}YYY@{}}
\toprule
\textbf{} & \textbf{Autoregressive LLM} & \textbf{Diffusion-style generator} \\
\midrule
Order & Left-to-right, one token at a time & Iterative denoise / unmask; can update many positions \\
Starts from & Prompt prefix & Noise / mask (formulation-dependent) \\
Phase-1 tools & KV cache, GQA, PagedAttention, speculation & Different bottleneck story \\
\bottomrule
\end{tabularx}
\end{table}


\textbf{Do not invent product claims} (``Model X ships diffusion language''). Stick to: AR is dominant for chat APIs; diffusion-for-language is an active alternative paradigm.


\bigskip


\clearpage
\section{Phase 3 --- Transformer architecture (interview depth)}


\textbf{Teaching video:} \href{https://www.youtube.com/watch?v=c533te7NSpI}{LLM Interview Series \#3}


Do \textbf{not} recite ``attention + FFN, done.'' Lead big picture $\rightarrow$ stack $\rightarrow$ one block $\rightarrow$ zoom attention.


\subsection{Step 3.1 --- Scope: decoder-only for chat}


Original ``Attention Is All You Need'' (Vaswani et al., 2017) had \textbf{encoder + decoder}. Modern chat models are \textbf{decoder-only}: generate the next token left-to-right. Unless asked for machine translation / seq2seq, focus the interview answer on the decoder stack.


\subsection{Step 3.2 --- Depth 1: whole LLM = three blocks}


\begin{table}[ht]
\centering\small
\begin{tabularx}{\textwidth}{@{}YY@{}}
\toprule
\textbf{Block} & \textbf{Role} \\
\midrule
\textbf{Input} & Tokens $\rightarrow$ token embeddings + positional info $\rightarrow$ input embeddings \\
\textbf{Processor} & Stacked Transformer blocks (this \textit{is} the architecture) \\
\textbf{Output} & Hidden state $\rightarrow$ logits over vocabulary $\rightarrow$ next-token choice \\
\bottomrule
\end{tabularx}
\end{table}


Place the Transformer in the \textbf{processor}. Ground it with a real prompt path: ``Give me a travel plan for Spain'' $\rightarrow$ embeddings $\rightarrow$ stack $\rightarrow$ logits $\rightarrow$ next token.


\textbf{Same diagram, two regimes:}

\begin{itemize}
\item Pretrain: embeddings / \(W_Q,W_K,W_V\) / FFN weights are \textbf{learned}.  
\item Inference: those weights are \textbf{fixed}; forward path shape is the same.
\end{itemize}


\subsection{Step 3.3 --- Depth 2: stack many identical blocks}


One block's output feeds the next: \(T_1 \to T_2 \to \cdots \to T_L\) (dozens to 100+ layers in large models). Composable stacking is a reason the architecture scaled.


Innovations cluster on two modules inside each block:

\begin{itemize}
\item \textbf{Attention} variants: MQA, GQA, MLA, \ldots{}  
\item \textbf{FFN} variants: MoE, SwiGLU, \ldots{}
\end{itemize}


\subsection{Step 3.4 --- Depth 3: modules inside one block}


Typical decoder block as taught in the lecture (exact norm placement varies --- pre-norm vs post-norm):


\begin{enumerate}
\item LayerNorm (or RMSNorm in many modern models)  
\item Multi-head attention  
\item Dropout  
\item Residual / skip  
\item LayerNorm / RMSNorm  
\item Feed-forward network (MLP)  
\item Dropout  
\item Residual / skip  
\end{enumerate}


\begin{table}[ht]
\centering\small
\begin{tabularx}{\textwidth}{@{}YY@{}}
\toprule
\textbf{Module} & \textbf{First-principles job} \\
\midrule
Norm (LN / RMSNorm) & Stabilize activations / training \\
Multi-head attention & \textbf{Only place tokens talk to each other} --- mix into context vectors. Causal mask: no future tokens \\
Dropout & Training regularization (classic DL) \\
Residual & Alternate path for gradients $\rightarrow$ fight vanishing gradients in deep stacks \\
FFN / MLP & \textbf{Per-token} expand (often \textasciitilde{}\(4\times d\)) $\rightarrow$ nonlinearity $\rightarrow$ contract. Tokens do \textbf{not} mix here. Reason: \textbf{expressivity} \\
\bottomrule
\end{tabularx}
\end{table}


\textbf{Critical contrast:} Attention = inter-token communication. FFN = intra-token wide transform. Without attention, ``Harry \ldots{} he'' never binds.


\subsection{Step 3.5 --- Depth 4: enough attention for this question}


\begin{enumerate}
\item \(X \rightarrow Q,K,V\) via \(W_Q,W_K,W_V\).  
\item Scores \(\propto QK\textasciicircum{}\top\); scale by \(\sqrt{d_k}\); causal mask; softmax $\rightarrow$ weights.  
\item Context \(\propto\) weights \(\times V\).  
\item Multi-head: split into heads; concat.
\end{enumerate}


Full KV-cache story is Interview \#1; GQA is Interview \#6.


\subsection{Step 3.6 --- Modern knobs (flag, don't invent)}


\begin{table}[ht]
\centering\small
\begin{tabularx}{\textwidth}{@{}YY@{}}
\toprule
\textbf{Older textbook} & \textbf{Common modern twist} \\
\midrule
Absolute positional embeddings & \textbf{RoPE} inside attention \\
LayerNorm & \textbf{RMSNorm} \\
GELU MLP & \textbf{SwiGLU} (gated) \\
Dense FFN & \textbf{MoE} routing (family-specific) \\
\bottomrule
\end{tabularx}
\end{table}


Mention as family-specific, not universal law.


\subsubsection{60-second script}

``Decoder-only LLM: embeddings in, stacked Transformer blocks, logits out. Each block is norm $\rightarrow$ causal MHA (tokens mix) $\rightarrow$ residual $\rightarrow$ norm $\rightarrow$ wide FFN (per-token expressivity) $\rightarrow$ residual. Attention is the only inter-token path; FFN lifts width then contracts. Variants like GQA/MoE change attention or FFN; serving then caches \(K,V\) at decode.''


\bigskip


\clearpage
\section{Phase 4 --- Attention in full (conceptual then matrix)}


\textbf{Teaching video:} \href{https://www.youtube.com/watch?v=cquX2tOODUI}{LLM Interview Series \#4}


\subsection{Step 4.1 --- Place it in the bird's-eye view first}


Input $\rightarrow$ Processor (attention sits inside each block after the first norm) $\rightarrow$ Output. Only then zoom.


\subsection{Step 4.2 --- Conceptual problem attention solves}


\paragraph{Setup.}

Language is full of long-range links:

\begin{itemize}
\item \texttt{Harry boarded the train. \textbf{He} \ldots{}} --- who is \textit{he}?  
\item \texttt{The dog chased the ball. \textbf{It} could not catch \textbf{it}.} --- first \textit{it}$\rightarrow$dog, second \textit{it}$\rightarrow$ball.  
\item Code: ``change the first 20 lines'' needs earlier context.
\end{itemize}


\paragraph{Problem.}

If every token were processed in isolation, the model would have \textbf{no clue about its neighbors}. Context would not exist.


\paragraph{Root cause.}

Next-token modeling needs \textbf{links} between a token and earlier tokens, plus \textbf{how much} to weight each earlier token.


\subsubsection{Solution in words}

Attention captures links: for a given query token, which past keys matter, and with what strength.


\textbf{One-liner:} Without attention, tokens are islands; with attention, each token can read a weighted view of its past.


\subsection{Step 4.3 --- Matrix mechanism (write on the board)}


Input embeddings \(X\) (e.g. 3 tokens $\times$ \(d\)). Trained in pretraining; frozen at inference.


\begin{enumerate}
\item \(Q = X W_Q,\quad K = X W_K,\quad V = X W_V\).  
\item Scores \(S = Q K\textasciicircum{}\top\). Row \(i\), column \(j\) = how much token \(i\) attends to token \(j\).  
\item \textbf{Causal mask:} upper triangle blocked --- no future tokens.  
\item Scale (typically \(/\sqrt{d_k}\)), \textbf{softmax} $\rightarrow$ attention weights.  
\item Context \(=\) weights \(\times V\) --- each position becomes a \textbf{context vector} mixed from values.
\end{enumerate}


Scores decide \textbf{where} to look; values carry \textbf{what} content is mixed.


\subsection{Step 4.4 --- Multi-head}


Single-head = one score matrix = one perspective.   Multi-head: multiple projection sets $\rightarrow$ multiple score matrices $\rightarrow$ concat.


\textbf{Why:} one sentence can need several bindings at once (\textit{artist painted \ldots{} woman with a brush} --- brush with artist vs brush with woman).


\subsubsection{60-second script}

``Attention is how tokens get context: without it they don't know their neighbors. For each position we form Q, K, V; scores are \(QK\textasciicircum{}\top\); causal mask blocks the future; softmax turns scores into weights; weighted values become the context vector. Multi-head repeats that with different projections.''


\bigskip


\clearpage
\section{Phase 5 --- Pretraining}


\subsection{Step 5.1 --- Data}


\paragraph{Setup.}

Pretraining is next-token prediction on massive text (web, books, code, \ldots{}).


\paragraph{Problem.}

Raw crawls contain spam, boilerplate, duplicates, PII. Duplicates inflate memorization and can contaminate evaluation later.


\paragraph{Root cause.}

Gradient updates follow the empirical data distribution. Garbage in $\rightarrow$ capability and safety issues out.


\paragraph{Solution.}

Filter, deduplicate, and \textbf{domain-mix} (web + code + books + \ldots{}), as described in modern open recipes (e.g. Llama-family reports). Synthetic rewritten data is sometimes used carefully; overuse risks distribution collapse.


\subsection{Step 5.2 --- Scaling laws}


\paragraph{Setup.}

Pretraining costs money. Choose model size and number of training tokens under a FLOP budget.


\paragraph{Problem.}

``Just make the model bigger'' can waste compute if the model is undertrained on too few tokens (or the reverse).


\paragraph{Root cause.}

Loss depends jointly on parameters, data, and compute --- not parameters alone.


\paragraph{Solution.}

Empirical scaling laws: Kaplan et al. (2020); Hoffmann et al. / Chinchilla (2022). Chinchilla emphasizes compute-optimal \textbf{tokens-per-parameter}. Fits are empirical; data quality shifts constants. This guides budgets; it does \textbf{not} guarantee every downstream skill scales smoothly.


\subsection{Step 5.3 --- What pretraining produces --- and what it does not}


After pretraining you have a \textbf{base model}: a strong next-token predictor.


It does \textbf{not} reliably:

\begin{itemize}
\item follow user instructions,  
\item refuse unsafe asks,  
\item use a chat format,  
\item prefer helpful answers among many fluent ones.
\end{itemize}


\textbf{Root cause:} the objective never saw an explicit ``be a helpful assistant'' reward --- only ``predict the next token on internet-like text.''


Hand the baton to \textbf{post-training}. Do not expect prompting alone to fully replace missing base competence.


\subsection{Step 5.4 --- Training systems (names after the problem)}


\paragraph{Setup.}

A frontier run does not fit on one GPU: weights, optimizer states, activations, long sequences.


\paragraph{Problem.}

Naive data-parallel replication of full Adam states OOMs; single-device batch is too small.


\subsubsection{Solution toolkit (standard, not a new paper claim)}

AdamW + schedules; mixed precision (BF16/FP16/FP8); gradient checkpointing; data / tensor / pipeline parallelism; ZeRO / FSDP; sequence/context parallelism for long context.


\bigskip


\clearpage
\section{Phase 6 --- Post-training}


\subsection{Step 6.1 --- Supervised Fine-Tuning (SFT)}


\paragraph{Setup.}

You have a base LM. Users want answers to instructions, not random web continuations.


\paragraph{Problem.}

Ask a base model a question; it may continue as if writing a webpage, ignore the ask, or produce the wrong format.


\paragraph{Root cause.}

Pretraining distribution $\neq$ assistant distribution.


\paragraph{Solution.}

Collect \((\mathrm{prompt}, \mathrm{desired\ response})\) demonstrations and fine-tune with ordinary likelihood. This is Step 1 of InstructGPT (Ouyang et al., 2022).


\paragraph{Payoff.}

Teaches format, tool schemas, basic instruction following. Still insufficient when many answers are fluent but differently preferred.


\subsection{Step 6.2 --- Classic RLHF (do not start with ``PPO'')}


\paragraph{Setup.}

After SFT, the model can produce multiple fluent answers. Humans still have preferences: more helpful, less toxic, more truthful.


\paragraph{Problem.}

SFT clones demonstrations. It does not directly optimize ``prefer A over B.'' Writing one gold answer for every open-ended prompt is incomplete.


\paragraph{Root cause.}

We need a ranking signal, then a way to optimize against it without reward hacking / drift.


\subsubsection{Solution --- InstructGPT three steps (Ouyang et al., 2022)}

\begin{enumerate}
\item \textbf{SFT} on demonstrations.  
\item \textbf{Reward model (RM):} humans compare outputs; train an RM to predict the preferred one.  
\item \textbf{PPO:} optimize the policy to increase RM reward, with a \textbf{KL penalty} toward a reference policy (limits over-optimization). The KL is the relative-entropy idea from Phase 2.
\end{enumerate}


\paragraph{Tradeoffs.}

Preference data cost, RM misspecification, reward hacking, multi-model training complexity (policy + value + RM + reference).


\textbf{Sources:} Ouyang et al., 2022; OpenAI instruction-following writeup.


\subsection{Step 6.3 --- DPO (problem first)}


\paragraph{Setup.}

RLHF works, but the stack is heavy: sample $\rightarrow$ RM $\rightarrow$ PPO with value model and KL control.


\paragraph{Problem.}

Engineering instability and cost make iteration slow for many labs.


\paragraph{Solution.}

\textbf{DPO} (Rafailov et al., 2023): train the policy directly on preference pairs with a classification-style objective that implicitly encodes the reward (under Bradley--Terry + KL-constrained optimal-policy assumptions). Related objectives (IPO, KTO, \ldots{}) exist for different data/failure modes --- name them only after stating which pain point they target.


\subsection{Step 6.4 --- Reasoning post-training (high level, no invented scores)}


Chat alignment $\neq$ contest math / coding reliability.


Families (keep claims checkable):

\begin{itemize}
\item \textbf{ORM} --- score final answer when a verifier exists.  
\item \textbf{PRM} --- score steps (denser credit).  
\item \textbf{RL with verifiable rewards} --- unit tests / exact checkers.  
\item \textbf{GRPO} (DeepSeekMath; used in DeepSeek-R1 line): group of outputs per prompt; group-normalized rewards as advantages; no separate value network.
\end{itemize}


Do \textbf{not} invent score tables.


\subsection{Step 6.5 --- Parameter-efficient fine-tuning}


Full copies of weights per task are impractical. Task updates are often low-rank relative to full \(W\). \textbf{LoRA} / \textbf{QLoRA} / adapters store small \(A,B\) instead. Fine-tune on smaller GPUs; serve multiple adapters on one base.


\bigskip


\clearpage
\section{Phase 7 --- RAG vs fine-tuning (application design)}


\textbf{Teaching video:} \href{https://www.youtube.com/watch?v=cCXjumE70-g}{LLM Interview Series \#8}


This is a \textbf{system-design} question. Interviewers test whether you can choose an approach for a vague industry request --- not recite two definitions.


\subsection{Step 7.1 --- The interview trap}


You may know RAG and fine-tuning in isolation. The hard part is connecting them: \textit{for this product, which one (or both), and why?}


Vague client: ``Build a chatbot for our domain.''   Strong answer: clarify \textbf{grounded lookup} of owned documents vs \textbf{new behavior / style / pattern} that is not literally a passage to retrieve.


\subsection{Step 7.2 --- What problem does RAG solve? (before the acronym)}


\paragraph{Setup.}

A nutrition / fitness company has \textasciitilde{}2,000 curated resources (their IP). Customers ask ``what should I eat for breakfast?'' Answers must be \textbf{grounded} in those resources --- not a generic web opinion.


\subsubsection{Naive solution}

Stuff \textbf{all} resources into the prompt every turn.


\paragraph{Problem.}

Corpus can be huge (teaching order-of-magnitude in the video: tens--hundreds of millions of tokens). Context windows are large but still limited (example contrast used in teaching: \textasciitilde{}1M $\ll$ 100M). You cannot fit the whole library.


\paragraph{Root cause.}

Context is a scarce resource (``like land''). Only a \textbf{slice} of the corpus is relevant per question.


\subsubsection{Solution --- Retrieval-Augmented Generation (Lewis et al., 2020)}

\begin{enumerate}
\item Chunk the corpus (paragraphs / sections --- a design choice).  
\item Embed chunks $\rightarrow$ vectors; embed the query.  
\item Retrieve top-\(k\) chunks by similarity (e.g. nearest neighbors / dot product).  
\item Put \textbf{query + retrieved chunks} into the LLM context; generate.
\end{enumerate}


Weights stay \textbf{fixed}. Name: \textbf{retrieval} augments \textbf{generation}.


\textbf{Production caveat from the teaching video:} embedding and storing vectors is expensive; stores often use quantization/compression. ``Vectorless RAG'' (tree-style) is mentioned as a fork --- do not invent details unless asked.


\textbf{Analogy:} open-book exam --- flip to relevant pages; book stays beside you.


\textbf{When you do not need RAG:} if owned knowledge fits in context with margin, pass it directly.


\subsection{Step 7.3 --- What is fine-tuning? (start from pretraining)}


\paragraph{Setup.}

A pretrained LM is billions of ``knobs'' fitted on internet-scale data. A raw pretrained checkpoint is often weak at following instructions (``convert 35 km to meters'').


\paragraph{Problem.}

You need narrower behavior (instructions, domain style, persona) without full pretraining again.


\paragraph{Solution.}

Curate a much smaller supervised set of (question, answer) pairs (teaching ballpark: order of \textasciitilde{}10k, not pretraining scale). Train again so weights \textbf{nudge}. That is fine-tuning.


\textbf{Key contrast:} fine-tuning \textbf{changes weights}; RAG does \textbf{not}.


\textbf{Pointer only:} the video links ``why small data works'' to intrinsic dimensionality of a pretrained model --- say ``small supervised nudge on a strong pretrained base,'' not a fake theorem.


\subsection{Step 7.4 --- Decision examples}


\subsubsection{Example A --- Digital clone $\rightarrow$ fine-tuning}

User dumps LinkedIn / YouTube / Instagram. You want a chatbot that \textbf{sounds like them} on \textit{new} questions not literally in the posts.


\begin{table}[ht]
\centering\small
\begin{tabularx}{\textwidth}{@{}YY@{}}
\toprule
\textbf{Approach} & \textbf{What happens} \\
\midrule
RAG & Retrieves relevant posts; does \textbf{not} teach speaking patterns for novel asks \\
Fine-tuning & Adjusts weights so patterns of speech are internalized \\
\bottomrule
\end{tabularx}
\end{table}


\subsubsection{Example B --- Pharma / nutrition grounded bot $\rightarrow$ start with RAG}

Answers must stay inside owned documents. Start with RAG: faster, cheaper, simpler. Caveat: RAG surfaces what's (near) in the store; it does not discover deep new patterns.


\subsubsection{Escalation}

\begin{lstlisting}
easiest $\rightarrow$ harder
RAG  $\rightarrow$  fine-tune pretrained model  $\rightarrow$  train SLM from scratch (pretrain)
\end{lstlisting}


\subsection{Step 7.5 --- Comparison table (draw this)}


\begin{table}[ht]
\centering\small
\begin{tabularx}{\textwidth}{@{}YYY@{}}
\toprule
\textbf{Axis} & \textbf{RAG} & \textbf{Fine-tuning} \\
\midrule
Finds patterns in data? & No --- retrieves existing text & Yes --- weights encode patterns \\
Cost & Lower & Higher \\
Simplicity & Higher & Lower \\
Personalization & Weak & Strong when trained for it \\
Analogy & Open-book exam & Study the night before (no book in the room) \\
Weights at serve & Unchanged & Changed (or LoRA/adapters) \\
\bottomrule
\end{tabularx}
\end{table}


Freshness: RAG wins when knowledge updates often (re-index). Many real systems do \textbf{both} --- say that \textit{after} the primary axes.


\subsubsection{60-second script}

``RAG and fine-tuning solve different failures. If a company's library won't fit in context, RAG retrieves relevant chunks and grounds the answer --- open-book; weights stay fixed. Fine-tuning nudges weights on a small labeled set so behavior or patterns change --- studying the night before. Grounded nutrition/pharma bot: start RAG. Digital clone on novel questions: fine-tune. Compare patterns, cost, simplicity, personalization; escalate RAG $\rightarrow$ FT $\rightarrow$ pretrain only as needed.''


\bigskip


\clearpage
\section{Phase 8 --- Inference engineering begins}


\textbf{Course framing:} \href{https://inference.vizuara.ai/}{inference.vizuara.ai}


\paragraph{Setup.}

Training builds weights. Users experience \textbf{serving}: tokens out per second, dollars per million tokens, GPU OOMs at 2 a.m.


\paragraph{Problem.}

A notebook \texttt{model.generate()} that feels fine for one prompt collapses under concurrent users, long contexts, and cost caps.


\paragraph{Root cause.}

Inference is autoregressive, often memory-bandwidth heavy, and multi-tenant. Training tricks do not automatically solve decode.


\paragraph{Discipline.}

Treat inference as its own stack: hardware $\rightarrow$ kernels $\rightarrow$ attention/KV design $\rightarrow$ scheduler $\rightarrow$ quantization $\rightarrow$ speculation.


\paragraph{Hardware lab mindset.}

The same 7B weights bottleneck differently on a laptop, Raspberry Pi, phone, or datacenter GPU. Measure tok/s, VRAM/RSS, quality per device. Do not quote an A100 blog post as universal truth.


\bigskip


\clearpage
\section{Phase 9 --- Prefill vs decode, latency, memory, roofline}


\subsection{Step 9.1 --- Prefill vs decode (must name this split)}


Example: prompt \texttt{A sunset is} $\rightarrow$ model continues \texttt{extremely beautiful\ldots{}}


\begin{enumerate}
\item \textbf{Prefill:} consume \textbf{all} prompt tokens (largely in parallel), run attention, emit the \textbf{first} output token. Often more compute-bound.  
\item \textbf{Decode:} emit \textbf{one new token at a time} until stop. This is where naive recompute hurts.
\end{enumerate}


Until first token $\rightarrow$ prefill; after that $\rightarrow$ decode. Every later optimization attaches to this split.


\textbf{Root cause of the split:} autoregressive factorization \(P(x_t\mid x_{<t})\) forces sequential new tokens even when the prompt can be processed in parallel.


\subsection{Step 9.2 --- Latency vs throughput}


\begin{itemize}
\item \textbf{Latency:} TTFT (time to first token), TPOT (time per output token) for \textit{one} user.  
\item \textbf{Throughput:} tokens/sec or requests/sec across \textit{many} users.
\end{itemize}


GPUs want large batches; users want snappy single-stream decode. Continuous batching (Phase 13) trades them consciously.


\subsection{Step 9.3 --- Memory bill of materials}


During inference, GPU/host memory typically holds:

\begin{enumerate}
\item \textbf{Weights}  
\item \textbf{Activations} (short-lived)  
\item \textbf{KV cache} (grows with sequence $\times$ layers $\times$ KV layout)
\end{enumerate}


At long context / high concurrency, \textbf{KV often dominates} and OOMs first.


\subsection{Step 9.4 --- Roofline intuition}


\paragraph{Setup.}

Every kernel needs FLOPs and bytes from HBM.


\paragraph{Problem.}

You ``optimize FLOPs'' but wall-clock does not move.


\paragraph{Root cause.}

Decode is often \textbf{memory-bandwidth bound} (reload weights + KV each step), not FLOP-bound.


\subsubsection{Solution direction}

Quantization, GQA/MQA, FlashAttention IO patterns, speculation --- attack \textbf{bytes moved} or \textbf{steps taken}.


\bigskip


\clearpage
\section{Phase 10 --- KV cache (full interview walkthrough)}


\textbf{Teaching video:} \href{https://www.youtube.com/watch?v=CxRGWfcGVbs}{LLM Interview Series \#1}


Do \textbf{not} open with ``KV cache makes generation faster.'' Construct from first principles.


\subsection{Step 10.1 --- It is an inference concept}

Shows up at serve / generate time. Weights \(W_Q,W_K,W_V\) are already fixed.


\subsection{Step 10.2 --- Prefill at matrix level}

Input embeddings \(X\) multiply fixed \(W_Q,W_K,W_V\) $\rightarrow$ \(Q,K,V\).   Scores \(\propto QK\textasciicircum{}\top\) (scale / softmax / causal) $\rightarrow$ weights. Context \(\propto\) weights \(\times V\). Last position $\rightarrow$ logits $\rightarrow$ first generated token.


\textbf{At end of prefill:} you already computed \(K\) and \(V\) rows for every prompt token. \textbf{Cache them.}


\subsection{Step 10.3 --- Naive decode (the wrong picture)}

To generate the next token, a naive approach recomputes \(Q,K,V\) for the \textbf{entire} sequence every step, full score matrix, full context matrix --- then keeps only the last row.


That recomputes past tokens' \(K,V\) that \textbf{cannot change} (same tokens, frozen weights).


\subsection{Step 10.4 --- Two realizations that define KV cache}


\textbf{Realization A.} Only the latest context vector is needed for the next token --- not a full context matrix of all positions for the logits step.


\textbf{Realization B.} That latest context still needs \textbf{full} past \(K\) and \textbf{full} past \(V\). The new query attends over all past keys; the weighted sum uses all past values.


So you need the whole \(K,V\) history --- but you must not \textbf{recompute} it.


\subsection{Step 10.5 --- Actual decode with cache}

\begin{enumerate}
\item Embed \textbf{only the new token} (\(1 \times d\)).  
\item Compute only that token's \(q_{\text{new}}, k_{\text{new}}, v_{\text{new}}\).  
\item \textbf{Append} \(k_{\text{new}}, v_{\text{new}}\) to cached \(K,V\).  
\item Use \(q_{\text{new}}\) against full \(K\textasciicircum{}\top\) $\rightarrow$ scores for this step only (\(1 \times\) seq).  
\item Mix with full \(V\) $\rightarrow$ one context vector $\rightarrow$ logits $\rightarrow$ next token.  
\item Repeat; cache grows by one \(K,V\) row per new token (per layer / head layout).
\end{enumerate}


\subsection{Step 10.6 --- Why not a Q cache?}

You need full \(K\), full \(V\), but \textbf{only} \(q\) for the \textbf{new} token. Past queries are not reused for the next step's scores. Hence \textbf{KV} cache, not QKV cache.


\subsection{Step 10.7 --- Payoff vs evil}


\begin{table}[ht]
\centering\small
\begin{tabularx}{\textwidth}{@{}YY@{}}
\toprule
\textbf{Good} & \textbf{Evil} \\
\midrule
Avoids redundant matmuls; compute closer to \textbf{linear} in new tokens vs \textbf{quadratic} recompute of the past & Memory grows with sequence $\times$ layers $\times$ KV heads $\times$ dim $\times$ 2 \\
Better interactive latency & Moving KV from HBM each step $\rightarrow$ bandwidth bottleneck $\rightarrow$ GQA, PagedAttention, quant \\
\bottomrule
\end{tabularx}
\end{table}


KV size scales roughly with:
\texttt{layers $\times$ kv\_heads $\times$ head\_dim $\times$ sequence\_length $\times$ bytes\_per\_element $\times$ 2}
(Exact formula depends on GQA/MQA and layout.)


\subsection{Step 10.8 --- Unknown output length (bridge)}

You do not know final decode length when the request arrives. Contiguous over-reservation $\rightarrow$ fragmentation $\rightarrow$ \textbf{PagedAttention} (Phase 12).


\subsubsection{Related names (state the problem first)}

\begin{table}[ht]
\centering\small
\begin{tabularx}{\textwidth}{@{}YYY@{}}
\toprule
\textbf{Idea} & \textbf{Problem} & \textbf{One line} \\
\midrule
Prefix / prompt caching & Repeated system prompts & Reuse KV for shared prefixes \\
Chunked prefill & Huge prefills block the GPU & Split prefill; interleave with decode \\
Eviction / compression (H2O, StreamingLLM, \ldots{}) & Infinite context, finite VRAM & Drop/compress less useful KV (method-specific) \\
\bottomrule
\end{tabularx}
\end{table}


\subsubsection{60-second script}

\begin{enumerate}
\item Inference-only; prefill vs decode.  
\item Prefill builds \(Q,K,V\); \textbf{store \(K,V\)}.  
\item Next token needs latest context $\Rightarrow$ full \(K,V\) + new \(q\) only.  
\item Don't recompute past \(K,V\); append new \(k,v\).  
\item No Q-cache.  
\item Speeds decode; costs memory/bandwidth.
\end{enumerate}


\bigskip


\clearpage
\section{Phase 11 --- Shrink the KV: MHA $\rightarrow$ MQA $\rightarrow$ GQA}


\textbf{Teaching video:} \href{https://www.youtube.com/watch?v=mtsY7JsGQjw}{LLM Interview Series \#6}


Shared setup: at decode, loading \(K,V\) from the cache is a major \textbf{HBM $\rightarrow$ compute} bandwidth cost.


\subsection{Step 11.1 --- Place attention, then multi-head}

LM = input $\rightarrow$ stacked Transformer blocks $\rightarrow$ logits. Inside a block: norm $\rightarrow$ \textbf{MHA} $\rightarrow$ residual $\rightarrow$ norm $\rightarrow$ FFN $\rightarrow$ residual.


Multi-head exists for \textbf{perspectives} (ambiguous attachment: brush with artist vs woman). Paper for MHA: Vaswani et al., 2017.


\subsection{Step 11.2 --- The need for GQA starts in the KV cache}

In MHA with \(h\) heads, decode stores per-head keys and values. \textbf{KV cache size scales with the number of KV heads.} Larger KV $\Rightarrow$ more bytes moved each decode step.


Innovation need: \textbf{shrink KV without throwing away too much quality.}


\subsection{Step 11.3 --- MQA (Shazeer, 2019) --- extreme fix}

Keep different queries per head, but force \textbf{one shared K and one shared V} across all heads.


\textbf{Store:} only one \(K\) and one \(V\) (not \(h\) of each) $\rightarrow$ large memory cut in the counting argument (up to \textasciitilde{}\(h\times\)).


\textbf{Why quality suffers:} heads diversify attention scores. If all heads share the same \(K\) (and \(V\)), diversity comes only from different \(Q\)s $\rightarrow$ weaker multi-perspective modeling. Teaching frame: MQA = low KV (good) + weaker language pattern capture (bad).


\subsection{Step 11.4 --- GQA (Ainslie et al., EMNLP 2023) --- middle}

Partition query heads into \(g\) groups. \textbf{Within a group}, share K and V. \textbf{Across groups}, K/V differ.


Toy with 4 query heads and 2 groups:

\begin{itemize}
\item Group 1: \(K_1=K_2\), \(V_1=V_2\)  
\item Group 2: \(K_3=K_4\), \(V_3=V_4\)  
\item But \(K_1 \neq K_3\)
\end{itemize}


Endpoints: \(g=h\) $\rightarrow$ MHA; \(g=1\) $\rightarrow$ MQA.


\begin{table}[ht]
\centering\small
\begin{tabularx}{\textwidth}{@{}YYY@{}}
\toprule
\textbf{} & \textbf{KV memory} & \textbf{Multi-perspective quality} \\
\midrule
\textbf{MHA} & Highest & Highest \\
\textbf{GQA} & Middle & Middle \\
\textbf{MQA} & Lowest & Lowest (often too weak) \\
\bottomrule
\end{tabularx}
\end{table}


Teaching practical claim: production open models rarely ship pure MQA; \textbf{GQA is the common compromise} (e.g. Llama family uses GQA; number of groups is a hyperparameter --- verify per model card, don't invent).


Honest drawback: GQA is middle for everything. MLA-style latent attention is a \textbf{different} mechanism --- do not equate it with GQA.


\subsubsection{60-second script}

Locate MHA $\rightarrow$ multiple heads for perspectives $\rightarrow$ MHA KV $\propto$ heads hurts decode $\rightarrow$ MQA shares one KV (extreme) $\rightarrow$ GQA shares KV per group (middle) $\rightarrow$ draw the spectrum.


\subsection{Step 11.5 --- Other long-context families (brief)}

\begin{table}[ht]
\centering\small
\begin{tabularx}{\textwidth}{@{}YYY@{}}
\toprule
\textbf{Family} & \textbf{Problem} & \textbf{Idea} \\
\midrule
Sliding window & Full \(n\textasciicircum{}2\) too costly & Attend inside a window \\
Sparse attention & Dense all-pairs wasteful & Structured sparsity \\
SSM / Mamba (Gu \& Dao et al., 2023) & Want better asymptotic cost & Selective state that summarizes the past \\
\bottomrule
\end{tabularx}
\end{table}


\bigskip


\clearpage
\section{Phase 12 --- FlashAttention (IO schedule, not new math)}


\subsubsection{Shared problem}

Standard attention materializes large \(S=QK\textasciicircum{}\top\) (and softmax) in HBM. Memory traffic dominates even when FLOPs look fine. Root cause: GPU \textbf{HBM $\leftrightarrow$ SRAM} IO.


\subsubsection{FlashAttention-1 (Dao et al., NeurIPS 2022)}

\begin{itemize}
\item Tiling: bring blocks of \(Q,K,V\) into SRAM.  
\item \textbf{Online softmax} so you never need the full score matrix in HBM.  
\item Recomputation in backward (training) to save memory.  
\item \textbf{Exact} attention (same math result, different IO schedule).
\end{itemize}


\subsubsection{FlashAttention-2}

Better work partitioning / parallelism; higher utilization; MQA/GQA support called out in FA2 materials. Reported \textasciitilde{}2$\times$ over FA1 in \textit{their} settings (hardware/workload dependent).


\subsubsection{FlashAttention-3 (Dao et al., NeurIPS 2024)}

Targets Hopper (H100) asynchrony (Tensor Cores + TMA): warp specialization, interleave matmul and softmax, FP8 path. Paper reports \textasciitilde{}1.5--2.0$\times$ vs FA2 on H100 in BF16/FP16 regimes --- \textbf{on that hardware/paper setting}.


\textbf{Interview line:} FlashAttention does not change mathematical attention; it changes the \textbf{IO schedule} so GPUs stop drowning in HBM traffic.


\bigskip


\clearpage
\section{Phase 13 --- Serving: continuous batching and PagedAttention}


\textbf{Teaching video:} \href{https://www.youtube.com/watch?v=-AB6m0Spo6c}{LLM Interview Series \#5}   \textbf{Paper/blog:} Kwon et al., SOSP 2023; \href{https://vllm.ai/blog/2023-06-20-vllm}{vLLM blog}


\subsection{Step 13.1 --- Why a serving engine $\neq$ \texttt{transformers.generate}}

HF generate is fine for demos. Multi-tenant production needs a scheduler, memory manager, and batching policy.


\subsection{Step 13.2 --- Continuous batching}


\paragraph{Setup.}

Requests arrive and finish at different times.


\paragraph{Problem.}

Static batching leaves GPU lanes idle when short sequences finish.


\paragraph{Solution.}

Dynamically add/remove sequences each iteration (iteration-level scheduling; Orca/vLLM lineage).


\subsection{Step 13.3 --- PagedAttention --- full layers (interview gold path)}


\subsubsection{Layer 1 --- Basics: GPU memory during inference}

\begin{enumerate}
\item Model weights (fixed after training)  
\item Activations (ephemeral)  
\item KV cache (grows; major during serving)
\end{enumerate}


\subsubsection{Layer 2 --- Problem: traditional KV allocation wastes memory}

You do \textbf{not} know final decode length in advance. Naive approach: reserve a \textbf{contiguous} KV slab per request up to \texttt{max\_tokens}.


Failures:

\begin{itemize}
\item \textbf{Reserved but unused} memory while a request is still generating (cannot give that RAM to another user).  
\item After a short request finishes early, free holes may be \textbf{too small / wrongly placed} for a waiting large request (\textbf{fragmentation}).
\end{itemize}


Result: low effective batch size $\rightarrow$ poor GPU utilization $\rightarrow$ low throughput.


vLLM authors report existing systems can waste a large fraction of KV memory via fragmentation and over-reservation (blog: on the order of \textbf{60--80%} waste in the regimes they studied). Treat as \textbf{their characterization}, not a universal constant.


\subsubsection{Layer 3 --- Root cause}

KV for a request is treated like one contiguous physical allocation. Variable lengths + contiguous reservation $\Rightarrow$ fragmentation.


\subsubsection{Layer 4 --- Solution: PagedAttention (+ vLLM)}

Inspired by OS \textbf{virtual memory / paging}:


\begin{enumerate}
\item Split KV memory into fixed-size \textbf{blocks} (pages). Block size is a system parameter (often discussed as \textbf{16 tokens} per block in vLLM materials).  
\item Each block stores K/V for that many tokens.  
\item Maintain a \textbf{free block list} and a \textbf{block table} mapping each request's logical token spans $\rightarrow$ physical blocks.  
\item Allocate blocks \textbf{on demand} as decode proceeds; a request's blocks need \textbf{not} be contiguous.  
\item Attention kernel gathers K/V by following the block table (\textbf{PagedAttention}).
\end{enumerate}


Additional capability from the same abstraction (\href{https://vllm.ai/blog/2023-06-20-vllm}{vLLM blog}): share blocks across sequences (e.g. parallel samples from one prompt) with reference counting / copy-on-write.


\subsubsection{Layer 5 --- Payoff}

Near-zero waste except partial last block (blog: under \textasciitilde{}\textbf{4%} waste in their characterization) $\rightarrow$ larger batches $\rightarrow$ higher throughput. Paper/blog report large gains vs prior systems \textbf{in their benchmarks}. \textbf{Do not treat a single speedup number as universal.}


\subsubsection{One engine step (mental model)}

Schedule requests $\rightarrow$ allocate/free KV blocks $\rightarrow$ run forward (prefill and/or decode) $\rightarrow$ sample $\rightarrow$ update caches / finish requests.


\bigskip


\clearpage
\section{Phase 14 --- Quantization}


\paragraph{Setup.}

Decode is often memory-bandwidth bound; weights + KV compete for VRAM.


\paragraph{Problem.}

FP16/BF16 models may not fit, or bandwidth dominates latency.


\subsubsection{Solution family (separate the names)}

\begin{itemize}
\item Precision ladder: FP16/BF16 $\rightarrow$ INT8 $\rightarrow$ INT4 (and related formats).  
\item Post-training methods: GPTQ, AWQ, GGUF ecosystems --- different recipes for \textit{which} tensors and \textit{how} calibrated.  
\item Can target weights, and sometimes KV/activations depending on the stack.
\end{itemize}


\paragraph{Interview discipline.}

Always say \textbf{which tensors}, \textbf{which bitwidth}, and \textbf{what quality/latency you measured}. Never claim ``INT4 is always better.''


\bigskip


\clearpage
\section{Phase 15 --- Speculative decoding}


\paragraph{Setup.}

Target-model decode is serial: one (or few) tokens per expensive forward.


\paragraph{Problem.}

Latency wall from sequential target steps.


\subsubsection{Solution --- draft--verify (Leviathan et al., 2023)}

A cheaper \textbf{draft} proposes several tokens; the \textbf{target} verifies them in parallel and accepts a prefix consistent with the target distribution (exact algorithms preserve the target's distribution when done correctly).


\subsubsection{Payoff / tradeoff}

Speedup depends on \textbf{acceptance rate}. Variants: small draft model, Medusa-style heads, EAGLE / EAGLE-3, MTP objectives. Name one method and its tradeoff; do not invent acceptance numbers.


\bigskip


\clearpage
\section{Phase 16 --- Sampling and frozen-weight procedures}


\textbf{Sampling:} temperature, top-\(k\), nucleus / top-\(p\) (Holtzman et al.), beam search. Changes behavior without retraining.


\textbf{With frozen weights:} few-shot / ICL (Brown et al., 2020); Chain-of-Thought (Wei et al., 2022); self-consistency (Wang et al., 2022); ReAct / tools; RAG (Phase 7); test-time compute (more samples / search / ``thinking'' tokens on hard items).


\bigskip


\clearpage
\section{Phase 17 --- Evaluation (brief but honest)}


\paragraph{Problem.}

A single leaderboard number can rise while users get worse or unsafe answers.


\paragraph{Root cause.}

Benchmarks measure narrow operationalized skills. Proxies can be gamed; contamination inflates scores; LLM-as-judge has biases.


\paragraph{Practice.}

Portfolio: capability suites, instruction rubrics, human preference (arena-style Elo), safety, \textbf{plus cost/latency}. Report threats honestly.


Routing systems (optional): report quality \textbf{and} cost (strong-call fraction, PGR) --- RouteLLM (Ong et al., ICLR 2025), Hybrid LLM (Ding et al., 2024).


\bigskip


\clearpage
\section{Phase 18 --- End-to-end in one breath}


\begin{lstlisting}
Text in the world
  $\rightarrow$ filter / dedup / mix + tokenize
  $\rightarrow$ pretrain next-token Transformer (+ distributed training)
  $\rightarrow$ base model
  $\rightarrow$ SFT $\rightarrow$ preferences (RLHF or DPO/\ldots{}) / reasoning RL
  $\rightarrow$ product: RAG and/or fine-tune
  $\rightarrow$ serve: prefill/decode, KV cache, GQA, FlashAttention,
           continuous batching + PagedAttention, quant, speculation
  $\rightarrow$ optional tools / test-time compute / routing
  $\rightarrow$ evaluate $\rightarrow$ iterate
\end{lstlisting}


\bigskip


\clearpage
\section{Phase 19 --- Explicit non-claims}


\begin{itemize}
\item No claim that one alignment method is universally best.  
\item No fabricated benchmarks or universal speedups.  
\item Preference alignment does not ``solve'' safety.  
\item Agents = orchestration + tools + memory \textbf{on top of} this stack.  
\item Hardware and traffic change which optimization wins --- measure.
\end{itemize}


\bigskip


\clearpage
\section{Phase 20 --- Landmark sources (only these; do not invent)}


\begin{table}[ht]
\centering\small
\begin{tabularx}{\textwidth}{@{}YY@{}}
\toprule
\textbf{Topic} & \textbf{Anchor} \\
\midrule
Transformer & Vaswani et al., 2017 \\
GPT-3 / few-shot & Brown et al., 2020 \\
Scaling & Kaplan et al., 2020; Hoffmann et al., 2022 \\
InstructGPT / RLHF & Ouyang et al., 2022 \\
DPO & Rafailov et al., 2023 \\
RAG & Lewis et al., 2020 \\
CoT / self-consistency & Wei et al., 2022; Wang et al., 2022 \\
Nucleus sampling & Holtzman et al. \\
MQA & Shazeer, 2019 \\
GQA & Ainslie et al., EMNLP 2023 \\
FlashAttention & Dao et al., 2022; FA3 NeurIPS 2024 \\
PagedAttention / vLLM & Kwon et al., SOSP 2023; vLLM blog 2023 \\
Speculative decoding & Leviathan et al., 2023 \\
Mamba & Gu \& Dao et al., 2023 \\
GRPO / R1 line & DeepSeekMath 2024; DeepSeek-R1 2025 \\
RouteLLM & Ong et al., ICLR 2025 \\
Teaching & Vizuara Interview \#1 KV, \#3 Transformer, \#4 Attention, \#5 PagedAttention, \#6 GQA, \#8 RAG vs FT; Mental Models \#5--\#7, \#9 \\
\bottomrule
\end{tabularx}
\end{table}


\bigskip


\clearpage
\section{Phase 21 --- Interview questions (with full answer skeletons)}


Speak \textbf{setup $\rightarrow$ problem $\rightarrow$ root cause $\rightarrow$ solution $\rightarrow$ payoff}. Prefer an example you thought through.


\subsection{Foundations}


\textbf{Q1. What is an LLM, really?}   Parameterized next-token model via \(P(x_t\mid x_{<t})\); train with CE; generate by sampling. Not a separate ``understanding'' module. (Brown et al. / standard LM.)


\textbf{Q2. Why Transformers instead of RNNs?}   Need parallel training and long-range mixing at scale; self-attention + FFN + causal mask. (Vaswani 2017.)


\textbf{Q3. Explain the Transformer for a chat model.}   Use Phase 3 script: three blocks $\rightarrow$ eight modules with \textit{why} $\rightarrow$ QKV zoom $\rightarrow$ variants live in attention/FFN. (\href{https://www.youtube.com/watch?v=c533te7NSpI}{Video \#3})


\textbf{Q4. What is attention?}   Use Phase 4 script: islands $\rightarrow$ Harry/he $\rightarrow$ QKV $\rightarrow$ causal $\rightarrow$ multi-head. (\href{https://www.youtube.com/watch?v=cquX2tOODUI}{Video \#4})


\textbf{Q5. What is entropy in LLMs?}   Expected surprise; CE trains; perplexity scores; temperature spends; KL limits RLHF drift. (\href{https://www.youtube.com/watch?v=q3agYqgqklU}{Video Mental Model \#9})


\textbf{Q6. Expressivity vs compression?}   Expressivity: lift when stuck (watch overfit). Compression: narrow waist, work small, project up.


\subsection{Pretraining / post-training}


\textbf{Q7. What does pretraining give / miss?}   Strong base next-token predictor; misses reliable instructions/preferences $\rightarrow$ post-training.


\textbf{Q8. What is SFT?}   Supervised \((\mathrm{prompt},\mathrm{response})\) fine-tune; InstructGPT step 1. (Ouyang 2022.)


\textbf{Q9. Explain RLHF without starting at PPO.}   Many fluent answers $\rightarrow$ preference signal $\rightarrow$ RM $\rightarrow$ optimize policy with KL to reference. Tradeoffs: cost, hacking, complexity.


\textbf{Q10. What problem does DPO solve?}   Heavy RM+PPO loop; train policy on preference pairs directly. (Rafailov 2023.)


\textbf{Q11. What is LoRA?}   Task updates often low-rank; store \(A,B\) with \(r\ll d\); cheaper multi-task fine-tunes.


\subsection{RAG vs fine-tuning}


\textbf{Q12. RAG vs fine-tuning --- when which?}   Use Phase 7 full story + table + escalation. (\href{https://www.youtube.com/watch?v=cCXjumE70-g}{Video \#8})


\textbf{Q13. Can you use both?}   Yes: FT for behavior/format; RAG for grounded facts --- after primary axes.


\subsection{Inference core}


\textbf{Q14. Prefill vs decode?}   Prefill = parallel over prompt, first token; decode = sequential; different bottlenecks.


\textbf{Q15. What is the KV cache? Why not cache Q?}   Use Phase 10 full walkthrough. (\href{https://www.youtube.com/watch?v=CxRGWfcGVbs}{Video \#1})


\textbf{Q16. Latency vs throughput?}   Name the product goal; continuous batching trades them.


\textbf{Q17. Why is decode often memory-bound?}   Each step reloads weights (+KV) from HBM for little new compute.


\subsection{Attention variants / kernels}


\textbf{Q18. MHA vs MQA vs GQA?}   Use Phase 11 spectrum story. (Ainslie 2023; \href{https://www.youtube.com/watch?v=mtsY7JsGQjw}{Video \#6})


\textbf{Q19. What is FlashAttention?}   Same math; IO-aware tiling + online softmax. (Dao et al.)


\subsection{Serving}


\textbf{Q20. What is PagedAttention?}   Use Phase 13 five layers exactly. (\href{https://www.youtube.com/watch?v=-AB6m0Spo6c}{Video \#5}; Kwon/vLLM)


\textbf{Q21. Continuous batching?}   Sequences finish at different times; dynamically admit/finish.


\textbf{Q22. Walk one vLLM-style engine step.}   Schedule $\rightarrow$ allocate/free KV blocks $\rightarrow$ forward $\rightarrow$ sample $\rightarrow$ update.


\subsection{Quant / speculation / design}


\textbf{Q23. Why quantize? Which tensors?}   Fit + bandwidth; specify tensors, bitwidth, method, measured quality.


\textbf{Q24. Speculative decoding?}   Draft proposes; target verifies; preserve target distribution; acceptance rate is the dial. (Leviathan 2023.)


\textbf{Q25. Design a low-latency, high-throughput serving system.}   Prefill/decode $\rightarrow$ KV good/evil $\rightarrow$ GQA $\rightarrow$ FlashAttention $\rightarrow$ continuous batching + PagedAttention $\rightarrow$ quant $\rightarrow$ optional speculation $\rightarrow$ measure TTFT, TPOT, throughput, VRAM; state tradeoffs.


\textbf{Q26. Domain chatbot tomorrow --- what first?}   Clarify grounded docs vs persona. Docs $\rightarrow$ RAG MVP. Style $\rightarrow$ FT/LoRA. Escalate with evidence.


\textbf{Q27. OOM at long context with many chats?}   KV size $\rightarrow$ GQA? PagedAttention/fragmentation? max-token reservation? quant? batch vs latency SLO?


\textbf{Q28. Accuracy fine, p95 latency bad?}   Prefill vs decode; TTFT vs TPOT; batching; KV bandwidth; speculation acceptance; FLOP vs bandwidth bound.


\subsection{Quick-fire openers}


\begin{table}[ht]
\centering\small
\begin{tabularx}{\textwidth}{@{}YY@{}}
\toprule
\textbf{If they ask\ldots{}} & \textbf{Open with\ldots{}} \\
\midrule
Entropy & Disorder/surprise $\rightarrow$ CE, temp, KL \\
Transformer & Input / processor / output, then modules \\
Attention & Tokens are islands $\rightarrow$ QKV $\rightarrow$ causal \\
KV cache & Prefill/decode recompute waste \\
GQA & MHA KV bandwidth $\rightarrow$ share KV in groups \\
FlashAttention & HBM traffic, not new math \\
PagedAttention & Contiguous KV fragmentation \\
Quantization & Memory-bound decode + which tensors \\
Speculative decoding & Serial target steps $\rightarrow$ draft--verify \\
RAG vs FT & Grounded retrieve vs weight patterns \\
\bottomrule
\end{tabularx}
\end{table}


\subsection{Universal answer contract}

\begin{enumerate}
\item Do \textbf{not} start with the definition of X.  
\item System context.  
\item Naive approach + break.  
\item Root design mistake.  
\item X + mechanism.  
\item Payoff + tradeoff.  
\item One credible source.
\end{enumerate}


\bigskip


\textit{Files: \texttt{research/llm\_concepts\_from\_scratch.md} (this) $\cdot$ \texttt{research/llm\_concepts\_from\_scratch.tex} (LaTeX twin)}


\begin{thebibliography}{99}

\bibitem{vaswani2017attention}
Vaswani et~al. Attention is all you need. NeurIPS 2017.

\bibitem{brown2020gpt3}
Brown et~al. Language models are few-shot learners. NeurIPS 2020.

\bibitem{kaplan2020scaling}
Kaplan et~al. Scaling laws for neural language models. 2020.

\bibitem{hoffmann2022chinchilla}
Hoffmann et~al. Training compute-optimal large language models (Chinchilla). 2022.

\bibitem{ouyang2022instructgpt}
Ouyang et~al. Training language models to follow instructions with human feedback. NeurIPS 2022.
\url{https://arxiv.org/abs/2203.02155}.

\bibitem{rafailov2023dpo}
Rafailov et~al. Direct Preference Optimization. NeurIPS 2023.

\bibitem{lewis2020rag}
Lewis et~al. Retrieval-Augmented Generation for Knowledge-Intensive NLP Tasks. NeurIPS 2020.

\bibitem{wei2022cot}
Wei et~al. Chain-of-Thought Prompting. NeurIPS 2022.

\bibitem{wang2022selfcons}
Wang et~al. Self-Consistency Improves Chain of Thought Reasoning. ICLR 2023.

\bibitem{holtzman2019nucleus}
Holtzman et~al. The Curious Case of Neural Text Degeneration. ICLR 2020.

\bibitem{shazeer2019mqa}
Noam Shazeer. Fast transformer decoding: One write-head is all you need. 2019.

\bibitem{ainslie2023gqa}
Ainslie et~al. GQA. EMNLP 2023.

\bibitem{dao2022flash}
Dao et~al. FlashAttention. NeurIPS 2022.

\bibitem{dao2024flash3}
Dao et~al. FlashAttention-3. NeurIPS 2024.

\bibitem{kwon2023vllm}
Kwon et~al. Efficient Memory Management for Large Language Model Serving with PagedAttention. SOSP 2023.

\bibitem{vllmblog2023}
Kwon \& Li. vLLM blog. \url{https://vllm.ai/blog/2023-06-20-vllm}.

\bibitem{leviathan2023speculative}
Leviathan et~al. Speculative decoding. ICML 2023.

\bibitem{gu2023mamba}
Gu \& Dao et~al. Mamba. 2023.

\bibitem{deepseekmath2024}
Shao et~al. DeepSeekMath. 2024.

\bibitem{deepseekr12025}
DeepSeek-AI. DeepSeek-R1. 2025.

\bibitem{ong2025routellm}
Ong et~al. RouteLLM. ICLR 2025.

\bibitem{vizuara-site}
Vizuara. Inference Engineering. \url{https://inference.vizuara.ai/}.

\bibitem{vizuara-entropy-yt}
Vizuara. Entropy Mental Model \#9. \url{https://www.youtube.com/watch?v=q3agYqgqklU}.

\bibitem{vizuara-diffusion-yt}
Vizuara. Diffusion Mental Model \#7. \url{https://www.youtube.com/watch?v=bTEfdB2D1Ek}.

\bibitem{vizuara-expressivity-yt}
Vizuara. Expressivity Mental Model \#6. \url{https://www.youtube.com/watch?v=YE0udJCgDqc}.

\bibitem{vizuara-compression-yt}
Vizuara. Compression Mental Model \#5. \url{https://www.youtube.com/watch?v=hGS6RbAYLl0}.

\bibitem{vizuara-kv-yt}
Vizuara. Interview \#1 KV Cache. \url{https://www.youtube.com/watch?v=CxRGWfcGVbs}.

\bibitem{vizuara-transformer-yt}
Vizuara. Interview \#3 Transformer. \url{https://www.youtube.com/watch?v=c533te7NSpI}.

\bibitem{vizuara-attention-yt}
Vizuara. Interview \#4 Attention. \url{https://www.youtube.com/watch?v=cquX2tOODUI}.

\bibitem{vizuara-paged-yt}
Vizuara. Interview \#5 PagedAttention. \url{https://www.youtube.com/watch?v=-AB6m0Spo6c}.

\bibitem{vizuara-gqa-yt}
Vizuara. Interview \#6 GQA. \url{https://www.youtube.com/watch?v=mtsY7JsGQjw}.

\bibitem{vizuara-rag-ft-yt}
Vizuara. Interview \#8 RAG vs Fine-Tuning. \url{https://www.youtube.com/watch?v=cCXjumE70-g}.

\end{thebibliography}
\end{document}
